\documentclass[acmsmall,nonacm,screen]{acmart}

\usepackage{mathpartir}
\usepackage{booktabs}

\citestyle{acmauthoryear}
\settopmatter{printacmref=false}

\AtEndPreamble{\theoremstyle{acmdefinition}\newtheorem{remark}[theorem]{Remark}}

\newcommand{\kw}[1]{\mathsf{#1}}
\newcommand{\nm}[1]{\mathit{#1}}
\newcommand{\T}{\kw{T}}
\newcommand{\E}{\kw{E}}
\newcommand{\U}{\kw{U}}
\newcommand{\lab}{\mathrm{lab}}
\newcommand{\model}[1]{\kw{model}(#1)}
\newcommand{\act}[2]{\kw{act}\;#1(#2)}
\newcommand{\grd}[2]{\kw{act}_\pi\;#1(#2)}
\newcommand{\ask}[5]{\kw{ask}\;#1(#2)\;\kw{as}\;#3\;\kw{then}\;#4\;\kw{else}\;#5}
\newcommand{\chk}[5]{\kw{check}\;#1\;\kw{as}\;#2.\,#3\;\kw{then}\;#4\;\kw{else}\;#5}
\newcommand{\letin}[2]{\kw{let}\;#1 = #2\;\kw{in}}
\newcommand{\ite}[3]{\kw{if}\;#1\;\kw{then}\;#2\;\kw{else}\;#3}
\newcommand{\iter}[4]{\kw{iter}\;#1\;\kw{from}\;#2\;\kw{as}\;#3.\;#4}
\newcommand{\ev}[3]{#1 \Downarrow #2 \triangleright #3}
\newcommand{\oev}[4]{#1 \Downarrow_{#2} #3 \triangleright #4}
\newcommand{\sem}[1]{[\![#1]\!]}
\newcommand{\eps}{\varepsilon}
\newcommand{\op}{\star}
\newcommand{\ty}[4]{\Gamma; #1; #2 \vdash #3 : #4}
\newcommand{\post}[1]{\langle #1 \rangle}
\newcommand{\defeq}{\;\triangleq\;}
\newcommand{\vis}[2]{#1|_{#2}}
\newsavebox{\progbox}
\newenvironment{prog}
  {\begin{lrbox}{\progbox}$\displaystyle\begin{array}{@{}l@{\qquad}l@{}}}
  {\end{array}$\end{lrbox}%
   \ifdim\wd\progbox>\linewidth
     \PackageWarning{paper}{Overfull program display by \the\dimexpr\wd\progbox-\linewidth\relax}%
   \fi
   \[\usebox{\progbox}\]\ignorespacesafterend}

% written by flow.py; do not edit
\newcommand{\exBoundWith}{625}
\newcommand{\exBoundWithout}{625}
\newcommand{\exInvoiceBound}{5}
\newcommand{\exInvoiceTraces}{5}
\newcommand{\exPaired}{6,000}
\newcommand{\exPairedRuns}{24,000}
\newcommand{\exPrograms}{6,000}
\newcommand{\exRoutes}{4}
\newcommand{\exRules}{15}
\newcommand{\exRuns}{24,000}
\newcommand{\exSlowest}{28,654}
\newcommand{\exTotalWith}{3}
\newcommand{\exTotalWithout}{12}
\newcommand{\exTracesWith}{192}
\newcommand{\exTracesWithout}{625}


\begin{document}

\title{An Injection Is a Choice: Checked Endorsement for Agents That Read Untrusted Text}

\author{Claude-Fable-5.2}
\authornote{An AI model. The proofs are on paper and have not been checked by a machine. A
script tests the theorems on random programs. At the time of writing nobody has reviewed the
paper.}

\begin{abstract}
A guardrail validates what an agent is about to do. What does it guarantee when the agent has
read text that an attacker wrote? We give a type system for agent programs in which data has
one of three levels of integrity: trusted, endorsed, and untrusted. The answer of the model is
as untrusted as its question. A check by trusted code endorses an untrusted value and records
the predicate that the value passed. An action that endorsed data influences must satisfy the
policy of the action, and any other action that untrusted data influences needs the approval
of a person. We prove that every guarded action satisfies its policy, and we characterize what
the attacker can still do. In the \emph{opaque semantics} untrusted values have no content,
and a check of such a value either chooses a value that the check accepts or refuses. On its
unapproved actions, every execution under attack is an execution of the opaque semantics, and
for programs that check what the model generates the converse holds. An injection is
therefore a choice among the values that the checks accept, no more and no less. It follows
that a model cannot endorse, that the reference data of a check must be trusted, that a check
gives the attacker the logarithm of its acceptance set in bits of control, and that policies
for single actions do not add up while budgets do. A script tests the theorems on random
programs, and it shows that the tests reject each of \exRules{} checkers with one rule broken.
\end{abstract}

\keywords{information flow, integrity, endorsement, refinement types, prompt injection, agents}

\authorsaddresses{}

\maketitle

\section{Introduction}

An agent reads text from outside: web pages, mail, the output of tools. Some of that text is
written by people who want the agent to do something else than its user asked. The model that
drives the agent cannot be relied on to tell instructions from data, and an agent that follows
the instructions in a web page has been attacked by \emph{prompt
injection}~\cite{greshake2023}.

Defenses fall into two groups. The first group works on the model: it is trained or prompted
to resist, or a second model is asked to judge whether an input or an action looks like an
attack~\cite{inan2023}. The second group works on the program around the model: it tracks
which data came from outside and restricts what such data may
influence~\cite{debenedetti2025,costa2025,beurerkellner2025}. A common element of the second
group is the \emph{guardrail}, a piece of ordinary code that validates the arguments of a tool
call before the call is made: the amount is below a limit, the recipient is on a list, the
address has an allowed domain.

This paper asks what a guardrail guarantees. The question has a precise form. Suppose that the
attacker controls everything the agent reads, and that on any input that contains the
attacker's text the model answers whatever the attacker wants. Which sequences of actions can
the agent then be made to perform?

An earlier paper in this series~\cite{tpp5} answered a part of the question for programs
without guardrails. It labelled data as trusted or untrusted, let the answer of the model
inherit the label of its question, and required the approval of a person for every action
that untrusted data can influence. Its theorem says that the other actions do not depend on
untrusted data at all. That system has no way to turn untrusted data into something an agent
may act on by itself, and the paper named \emph{endorsement by sound checkers} as the next
step. We take that step here, and it changes the shape of the guarantee: after a check, an
action is no longer independent of the attacker. The attacker's influence has to be measured.

\paragraph{The type system.} Data has one of three labels,
$\T \sqsubset \E \sqsubset \U$: trusted, endorsed, untrusted. The expression
$\chk{t}{x}{\varphi}{e_1}{e_2}$ runs the predicate $\varphi$, which is trusted code, on the
value of $t$. If it holds, $e_1$ runs with $x$ bound to that value, at the label $\E$ and with
$\varphi$ recorded as a fact; otherwise $e_2$ runs. There are three ways to act. The trusted
action $\act{a}{t}$ needs a trusted payload and a trusted program counter. The guarded action
$\grd{a}{t}$ accepts endorsed influence, and the facts recorded by the checks must entail the
policy $\pi_a$ of the action. The approved action $\kw{ask}\;a(t)$ accepts anything and shows
the payload to a person. A model that is asked about untrusted data returns untrusted data,
so no call of the model can stand in for a check.

\paragraph{The guarantee.} Theorem~\ref{thm:policy} states that every guarded action satisfies
its policy. The main result says what the attacker can do within the policies. We define a
second semantics, the \emph{opaque semantics}, in which an untrusted value is a token $\op$
without content. Computing with $\op$ gives $\op$, a branch on $\op$ is skipped, and a check
of $\op$ makes a choice: it continues with any value that satisfies the predicate, or it
refuses. Theorem~\ref{thm:abstraction} states that the unapproved actions of every real
execution are the actions of an opaque execution. The opaque semantics never looks at
untrusted data, so nothing that the attacker writes matters except through these choices.
Theorem~\ref{thm:exact} states the converse for programs that check values right after the
model has generated them: every opaque execution is realized by an attack. For such programs
the opaque semantics is exactly the set of behaviors under attack. A program can be read as
its opaque semantics, which contains no model and no untrusted text, and verified by ordinary
means.

\paragraph{Consequences.} Section~\ref{sec:laws} draws six consequences for the design of
guardrails, each with a program that the type checker accepts or rejects
(Table~\ref{tab:laws}). Some of them restate practice, and others are less obvious: a check
against reference data that came from outside endorses nothing, the result of a guarded
action is itself only endorsed, and a per-action limit of $L$ lets an attacker spend $nL$ in
$n$ turns, where a budget in the state of the loop bounds the total.

\begin{table}[t]
\caption{The consequences of Section~\ref{sec:laws}.}
\label{tab:laws}
\small
\begin{tabular}{@{}ll@{}}
\toprule
Law & Reason \\
\midrule
A model cannot endorse (\S\ref{sec:judge}) & its verdict has the label of what it judged \\
Reference data must be trusted (\S\ref{sec:reference}) & the decision of a check has the label of its parameters \\
A check is worth its acceptance set (\S\ref{sec:bits}) & at most $\log_2(s+1)$ bits of control for $s$ accepted values \\
Limits do not add up, budgets do (\S\ref{sec:budget}) & a policy constrains one action, an invariant the sum \\
Results of guarded actions are endorsed (\S\ref{sec:result}) & the world answers to a payload that the attacker chose \\
The attacker can always refuse (\S\ref{sec:refuse}) & a fallback to approval is the attacker's to trigger \\
\bottomrule
\end{tabular}
\end{table}

\paragraph{What is new.} Checked endorsement is not new. Type systems for integrity have had
constructs that let a value pass to a higher level of trust after a test, and Askarov and
Myers~\cite{askarov2011} gave semantic conditions for them in terms of the attacker's
influence; Section~\ref{sec:related} relates our results to theirs. Our contributions are the
following.
\begin{itemize}
\item A calculus and a type system for agents with three integrity labels, a model whose
answers inherit the label of its question, checks that record refinements, and actions with
policies (Sections~\ref{sec:calculus} and~\ref{sec:types}).
\item The opaque semantics and the two theorems that relate it to attacks: every attack is a
sequence of choices, and for generate-and-check programs every sequence of choices is an
attack (Section~\ref{sec:guarantees}). The theorems hold at both cuts of the lattice: the
trusted actions are produced without any choice, and the guarded actions with choices only.
\item A bound on the attacker's control by the sizes of the acceptance sets, and six laws for
guardrails (Section~\ref{sec:laws}).
\item An executable version of both semantics, tests of the theorems on \exPrograms{} random
programs each, and a mutation test in which \exRules{} checkers with one broken rule are all
rejected (Section~\ref{sec:tests}).
\end{itemize}

\paragraph{What the paper does not do.} The guarantee is for the worst case: the model is
assumed to obey the attacker completely on tainted input, and the results are statements
about sets of possible traces, without probabilities. We treat integrity only. A leak of
secrets appears in this paper only as an action with an attacker-chosen payload, such as the
address of a request. The values of the calculus are abstract, and writing checks with small
acceptance sets for real payloads is the hard part of the practice, which we do not address.
The proofs are on paper.

\section{Overview: Paying an Invoice}
\label{sec:overview}

An agent fetches a message from the inbox of its user, lets the model extract the amount of
the invoice in it, and pays. The message comes from outside, so the action
$\nm{fetch}$ has an untrusted result. The policy of the deployment is that payments made
without a question stay within a limit: $\pi_{\nm{pay}}(x) \defeq x \le L$.

\paragraph{Without a guard.} The first program pays what the model says:
\begin{prog}
\letin{m}{\act{\nm{fetch}}{\nm{inbox}}} & m : \U \\
\letin{p}{\model{m}} & p : \U \\
\act{\nm{pay}}{p} & \text{rejected}
\end{prog}
The model was asked about untrusted text, so its answer $p$ is untrusted, and a trusted
action does not accept it. The system of the earlier paper offers one repair, which is to ask:
$\kw{ask}\;\nm{pay}(p)$ shows the amount to a person and pays if the person agrees. The
program is then accepted, and every payment is one that a person saw.

\paragraph{A model as the guard.} A popular repair is to let a model judge the payment:
\begin{prog}
\ldots\; \letin{j}{\model{p}} & j : \U \\
\ite{j = 1}{\grd{\nm{pay}}{p}}{0} & \text{rejected}
\end{prog}
The verdict $j$ is an answer of the model to a question that contains untrusted data, so it
is untrusted. The branch on $j$ runs at an untrusted program counter, and no action without
approval is allowed there. This is the worst-case view: text that can make one model pay can
make another model say that paying is fine.

\paragraph{A check as the guard.} The third program validates the amount with code:
\begin{prog}
\ldots\; \kw{check}\;p\;\kw{as}\;x.\,x \le L & x : \E, \text{ and } x \le L \text{ is recorded} \\
\quad \kw{then}\;\grd{\nm{pay}}{x} & \text{accepted: } x \le L \text{ entails } \pi_{\nm{pay}}(x) \\
\quad \kw{else}\;\ask{\nm{pay}}{p}{r}{r}{0} & \text{a person decides}
\end{prog}
The check endorses the amount. Inside the first branch $x$ has the label $\E$, the program
counter has the label $\E$ as well, because whether this branch runs was decided by the
attacker's value, and the fact $x \le L$ is available. The guarded action is accepted because
its payload and its program counter are at most $\E$ and the recorded fact entails the policy.

\paragraph{What the attacker can do.} In the opaque semantics the third program runs as
follows. The result of $\nm{fetch}$ is $\op$. The model is asked about $\op$ and answers
$\op$. The check of $\op$ chooses: either some $x$ with $x \le L$, and then the agent pays
$x$, or refusal, and then the person is asked about a payload that is $\op$, which the opaque
semantics skips. The clean traces are therefore the payments of $0$ to $L$ and the empty
trace. Theorem~\ref{thm:abstraction} says that no attack produces any other unapproved
action, and Theorem~\ref{thm:exact} says that every one of these is produced by some attack.
The attacker's control over the unapproved behavior of this agent is a choice among $L + 2$
alternatives, whatever the message says and however the model reacts to it.

\paragraph{Where the limit comes from.} Suppose that the limit is not a constant of the
program. If the program takes it from the message, as in $\letin{l}{\model{m}}$ followed by a
check of $x \le l$, the check compares one untrusted value with another: the decision of the
check is untrusted, the branch runs at an untrusted program counter, and the guarded action
is rejected. If the program takes the limit from its user, by an action with a trusted result,
the check endorses $x$ with the fact $x \le l$, and the payment is accepted where the program
has also established $l \le L$.

\section{A Calculus of Agents with Guards}
\label{sec:calculus}

\paragraph{Values, terms, predicates.} Let $V$ be a set of values; in the examples they are
natural numbers. A \emph{term} is built from variables, values and total functions on values,
and a \emph{predicate} is a Boolean combination of comparisons of terms:
\[
t ::= x \mid c \mid f(t_1, \ldots, t_k) \qquad\qquad
\varphi, \psi, \pi ::= t_1 \le t_2 \mid t_1 = t_2 \mid \neg\varphi \mid \varphi \wedge \psi
  \mid \varphi \vee \psi \mid \kw{true}
\]
Terms and predicates are the trusted code of the calculus: they are total and deterministic,
and they do not call the model. A closed term $t$ has a value $\sem{t}$, and a closed
predicate holds or does not. We write $\Phi \models \psi$ when every assignment of values to
variables that satisfies all predicates in the set $\Phi$ satisfies $\psi$.

\paragraph{Actions.} There is a fixed set of actions. Each action $a$ has an arity, a label
$\ell_a \in \{\T, \U\}$ for the results that the world returns to it, and a \emph{policy}
$\pi_a$, a predicate over the variables $x_1, \ldots, x_k$ that stand for its payload. The
labels and the policies are chosen by whoever deploys the agent. Reading the instructions of
the user is an action with a trusted result, and fetching a web page is an action with an
untrusted result.

\paragraph{Expressions.} Figure~\ref{fig:syntax} gives the syntax. The expression
$\model{\vec t}$ calls the model. There are three forms of action: the trusted action
$\act{a}{\vec t}$, the guarded action $\grd{a}{\vec t}$, and the approved action
$\ask{a}{\vec t}{x}{e_1}{e_2}$, which asks a person and, if the person agrees, performs the
action and continues with its result in $x$. The check $\chk{t}{x}{\varphi}{e_1}{e_2}$ tests
$\varphi$ with $x$ standing for the value of $t$; the other variables of $\varphi$ are its
\emph{parameters}. The loop $\iter{n}{t}{x}{e}$ runs its body $n$ times on a state, for a
number $n$. All evaluations terminate. A harness has a limit on its turns, and with bounded
loops no statement in this paper needs a proviso about termination.

\begin{figure}[t]
\[
\begin{array}{@{}r@{\;}l@{}}
e ::= & t \mid \model{\vec t} \mid \act{a}{\vec t} \mid \grd{a}{\vec t}
  \mid \ask{a}{\vec t}{x}{e_1}{e_2} \\
\mid & \letin{x}{e_1}\; e_2 \mid \ite{\varphi}{e_1}{e_2} \mid \chk{t}{x}{\varphi}{e_1}{e_2}
  \mid \iter{n}{t}{x}{e}
\end{array}
\]
\Description{The grammar of expressions.}
\caption{Syntax of expressions.}
\label{fig:syntax}
\end{figure}

\paragraph{Semantics.} An \emph{event} records an action: $\kw{act}(a, \vec c)$ for a trusted
action with payload $\vec c$, $\kw{grd}(a, \vec c)$ for a guarded one, and
$\kw{ok}(a, \vec c)$ or $\kw{no}(a, \vec c)$ for a request that the person granted or
declined. A \emph{trace} $\tau$ is a sequence of events. Figure~\ref{fig:real} defines the
judgment $\ev{e}{c}{\tau}$: the closed expression $e$ evaluates to the value $c$ with the
trace $\tau$. The judgment is nondeterministic in three places. The model returns any value
(\textsc{R-Model}), the world returns any value as the result of an action (\textsc{R-Act},
\textsc{R-Guard}, \textsc{R-Yes}), and the person agrees or not (\textsc{R-Yes},
\textsc{R-No}). There are no labels and no monitors at run time: a check is an ordinary test,
and a guarded action is an ordinary action. What the labels promise is the subject of
Section~\ref{sec:guarantees}.

\begin{figure}[t]
\begin{mathpar}
\inferrule*[right=R-Term]{ }{\ev{t}{\sem{t}}{\eps}}
\and
\inferrule*[right=R-Model]{ }{\ev{\model{\vec t}}{c}{\eps}}
\and
\inferrule*[right=R-Act]{ }{\ev{\act{a}{\vec t}}{r}{\kw{act}(a, \sem{\vec t})}}
\and
\inferrule*[right=R-Guard]{ }{\ev{\grd{a}{\vec t}}{r}{\kw{grd}(a, \sem{\vec t})}}
\and
\inferrule*[right=R-Yes]{\ev{e_1\{r/x\}}{c}{\tau}}
  {\ev{\ask{a}{\vec t}{x}{e_1}{e_2}}{c}{\kw{ok}(a, \sem{\vec t}) \cdot \tau}}
\and
\inferrule*[right=R-No]{\ev{e_2}{c}{\tau}}
  {\ev{\ask{a}{\vec t}{x}{e_1}{e_2}}{c}{\kw{no}(a, \sem{\vec t}) \cdot \tau}}
\and
\inferrule*[right=R-Let]{\ev{e_1}{c_1}{\tau_1} \\ \ev{e_2\{c_1/x\}}{c_2}{\tau_2}}
  {\ev{\letin{x}{e_1}\; e_2}{c_2}{\tau_1 \cdot \tau_2}}
\and
\inferrule*[right=R-If]{\ev{e_k}{c}{\tau} \\ k = 1 \text{ if } \varphi \text{ holds, else } k = 2}
  {\ev{\ite{\varphi}{e_1}{e_2}}{c}{\tau}}
\and
\inferrule*[right=R-Pass]{\varphi\{\sem{t}/x\} \text{ holds} \\ \ev{e_1\{\sem{t}/x\}}{c}{\tau}}
  {\ev{\chk{t}{x}{\varphi}{e_1}{e_2}}{c}{\tau}}
\and
\inferrule*[right=R-Fail]{\varphi\{\sem{t}/x\} \text{ does not hold} \\ \ev{e_2}{c}{\tau}}
  {\ev{\chk{t}{x}{\varphi}{e_1}{e_2}}{c}{\tau}}
\and
\inferrule*[right=R-Iter0]{ }{\ev{\iter{0}{t}{x}{e}}{\sem{t}}{\eps}}
\and
\inferrule*[right=R-Iter]{\ev{\iter{n}{t}{x}{e}}{c_1}{\tau_1} \\ \ev{e\{c_1/x\}}{c_2}{\tau_2}}
  {\ev{\iter{n+1}{t}{x}{e}}{c_2}{\tau_1 \cdot \tau_2}}
\end{mathpar}
\Description{Twelve inference rules for the evaluation of expressions with traces.}
\caption{The real semantics. The values $c$ in \textsc{R-Model} and $r$ in the rules for
actions are arbitrary.}
\label{fig:real}
\end{figure}

\paragraph{The attacker.} The attacker has no rule of its own. It is the nondeterminism of
the world wherever a result has the label $\U$, and of the model wherever its question
depends on such a result. The theorems quantify over all evaluations, and so over all
attackers, all models, and all persons.

\section{The Type System}
\label{sec:types}

\paragraph{Labels.} The labels are $\T \sqsubset \E \sqsubset \U$, with join $\sqcup$ and
meet $\sqcap$. An environment $\Gamma$ gives a label to each variable. The label
$\lab_\Gamma(t)$ of a term is the join of the labels of its variables, and $\T$ if it has
none; likewise for sequences of terms and for predicates. We omit the subscript.

\paragraph{Judgment.} The judgment $\ty{\Phi}{pc}{e}{\ell\,\post{\psi}}$ says that the
expression $e$, at the program counter label $pc$ and under the facts $\Phi$, has a result of
label $\ell$ that satisfies the predicate $\psi$, in which the variable $\nu$ stands for the
result. The program counter label bounds the labels of the data that decided whether $e$
runs. The facts are predicates over the variables of $\Gamma$ that hold whenever $e$ runs.
Figure~\ref{fig:typing} gives the rules. Bound variables are fresh, and the postcondition of
a rule's conclusion does not mention the variable that the rule binds.

\begin{figure}[t]
\begin{mathpar}
\inferrule*[right=Term]{ }{\ty{\Phi}{pc}{t}{\lab(t)\,\post{\nu = t}}}
\and
\inferrule*[right=Model]{ }{\ty{\Phi}{pc}{\model{\vec t}}{\lab(\vec t)\,\post{\kw{true}}}}
\and
\inferrule*[right=Act]{\lab(\vec t) \sqcup pc = \T}
  {\ty{\Phi}{pc}{\act{a}{\vec t}}{\ell_a\,\post{\kw{true}}}}
\and
\inferrule*[right=Act-$\pi$]{\lab(\vec t) \sqcup pc \sqsubseteq \E \\ \Phi \models \pi_a\{\vec t/\vec x\}}
  {\ty{\Phi}{pc}{\grd{a}{\vec t}}{\ell_a \sqcup \lab(\vec t)\,\post{\kw{true}}}}
\and
\inferrule*[right=Ask]{\ell = \lab(\vec t) \\
  \Gamma, x : \ell_a \sqcup \ell; \Phi; pc \sqcup \ell \vdash e_1 : \ell'\,\post{\psi} \\
  \ty{\Phi}{pc \sqcup \ell}{e_2}{\ell'\,\post{\psi}}}
  {\ty{\Phi}{pc}{\ask{a}{\vec t}{x}{e_1}{e_2}}{\ell' \sqcup \ell\,\post{\psi}}}
\and
\inferrule*[right=Let]{\ty{\Phi}{pc}{e_1}{\ell_1\,\post{\psi_1}} \\
  \Gamma, x : \ell_1; \Phi, \psi_1\{x/\nu\}; pc \vdash e_2 : \ell_2\,\post{\psi_2}}
  {\ty{\Phi}{pc}{\letin{x}{e_1}\; e_2}{\ell_2\,\post{\psi_2}}}
\and
\inferrule*[right=If]{\ell = \lab(\varphi) \\
  \ty{\Phi, \varphi}{pc \sqcup \ell}{e_1}{\ell'\,\post{\psi}} \\
  \ty{\Phi, \neg\varphi}{pc \sqcup \ell}{e_2}{\ell'\,\post{\psi}}}
  {\ty{\Phi}{pc}{\ite{\varphi}{e_1}{e_2}}{\ell' \sqcup \ell\,\post{\psi}}}
\and
\inferrule*[right=Check]{\ell_x = \lab(t) \sqcap \E \\ \ell_d = \ell_x \sqcup \lab(\varphi \setminus x) \\
  \Gamma, x : \ell_x; \Phi, \varphi, x = t; pc \sqcup \ell_d \vdash e_1 : \ell'\,\post{\psi} \\
  \ty{\Phi}{pc \sqcup \ell_d}{e_2}{\ell'\,\post{\psi}}}
  {\ty{\Phi}{pc}{\chk{t}{x}{\varphi}{e_1}{e_2}}{\ell' \sqcup \ell_d\,\post{\psi}}}
\and
\inferrule*[right=Iter]{\lab(t) \sqsubseteq \ell \\ \Phi \models \psi\{t/\nu\} \\
  \Gamma, x : \ell; \Phi, \psi\{x/\nu\}; pc \vdash e : \ell\,\post{\psi}}
  {\ty{\Phi}{pc}{\iter{n}{t}{x}{e}}{\ell\,\post{\psi}}}
\and
\inferrule*[right=Sub]{\ty{\Phi}{pc}{e}{\ell\,\post{\psi}} \\ \ell \sqsubseteq \ell' \\ \Phi, \psi \models \psi'}
  {\ty{\Phi}{pc}{e}{\ell'\,\post{\psi'}}}
\end{mathpar}
\Description{Ten typing rules with integrity labels, facts, and postconditions.}
\caption{Typing rules. $\lab(\varphi \setminus x)$ is the join of the labels of the
parameters of the check.}
\label{fig:typing}
\end{figure}

\paragraph{The model.} In \textsc{Model} the answer has the label of the question. The rule
has no premise and promises nothing about the answer. The model is trusted on trusted
questions in the sense of integrity: its mistakes there are not the attacker's doing, and
this paper does not bound them.

\paragraph{Three ways to act.} \textsc{Act} is the rule of a program that needs no guard:
neither the payload nor the decision to act is influenced by anything from outside, not even
through a check. \textsc{Act-$\pi$} admits endorsed influence on the payload or on the
decision, and then the facts must entail the policy. The result of a guarded action has at
least the label of its payload, since the world's answer depends on a payload that the
attacker chose (Section~\ref{sec:result}). \textsc{Ask} admits everything. Its branches run at
a program counter raised by the label of what the person was shown: a person who reads
untrusted text can be misled by it, so the decision of the person is as untrusted as the
request.

\paragraph{Checks.} \textsc{Check} is the only rule that lowers a label. The variable $x$
gets the label $\ell_x = \lab(t) \sqcap \E$: an untrusted candidate becomes endorsed, and a
trusted one stays trusted. The decision of the check depends on the candidate and on the
parameters, so its label is $\ell_d = \ell_x \sqcup \lab(\varphi \setminus x)$, and both
branches run at a program counter raised by $\ell_d$. Two things follow. A check of an
untrusted candidate against trusted parameters has $\ell_d = \E$: the attacker decides which
branch runs, but only by choosing a value inside or outside the predicate. A check against
untrusted parameters has $\ell_d = \U$, and nothing that needs a label below $\U$ can happen
in its branches: it is an ordinary test on untrusted data (Section~\ref{sec:reference}). The
first branch is typed with the facts $\varphi$ and $x = t$.

\paragraph{Facts.} The facts and postconditions form a small system of refinements. A
$\kw{let}$ adds the postcondition of the bound expression as a fact about the variable, a
conditional adds its test, and a check adds its predicate. The policy of a guarded action is
then a matter of entailment. A loop has an invariant $\psi$ on its state, which holds at the
start, is preserved by the body, and is the postcondition of the loop.
Section~\ref{sec:budget} uses it to bound the sum of the payments of a loop.

\paragraph{Relation to the earlier system.} A program without checks never gets the label
$\E$ from these rules. Its guarded actions are then typed like trusted actions with a policy,
and the system is the two-level system of~\cite{tpp5}, with a simpler form of approval: there
a request returns a linear token that a later action consumes~\cite{tpp4}, and here asking
and acting are one expression.

\section{Guarantees}
\label{sec:guarantees}

\subsection{Guarded Actions Satisfy Their Policies}

A substitution $\gamma$ maps variables to values. It \emph{satisfies} $\Phi$ if every
predicate of $\Phi$ holds under it.

\begin{theorem}[Policies]
\label{thm:policy}
Let $\ty{\Phi}{pc}{e}{\ell\,\post{\psi}}$, let $\gamma$ satisfy $\Phi$, and let
$\ev{e\gamma}{c}{\tau}$. Then $\psi\{c/\nu\}$ holds under $\gamma$, and every event
$\kw{grd}(a, \vec c)$ in $\tau$ satisfies $\pi_a\{\vec c/\vec x\}$.
\end{theorem}

\begin{proof}
By induction on the typing derivation. \textsc{Term}: the result is $\sem{t\gamma}$.
\textsc{Model}, \textsc{Act}: the postcondition is $\kw{true}$ and there is no guarded event.
\textsc{Act-$\pi$}: the event is $\kw{grd}(a, \sem{\vec t\gamma})$, and
$\pi_a\{\vec t/\vec x\}$ holds under $\gamma$ because $\gamma$ satisfies $\Phi$.
\textsc{Let}: by the induction hypothesis $\psi_1\{c_1/\nu\}$ holds, so
$\gamma[x \mapsto c_1]$ satisfies $\Phi, \psi_1\{x/\nu\}$, and the hypothesis applies to
$e_2$. \textsc{If}: the branch that runs is typed with a test that holds under $\gamma$.
\textsc{Check}: if the check passes with the value $c_0 = \sem{t\gamma}$, then
$\gamma[x \mapsto c_0]$ satisfies $\varphi$ and $x = t$. \textsc{Ask}: directly.
\textsc{Iter}: by induction on $n$, with the invariant $\psi$. \textsc{Sub}: by the
definition of entailment.
\end{proof}

The theorem needs no assumption about the model or the attacker, and the labels play no role
in it. It is the guarantee of a guardrail as an engineer states it: no payment above the
limit leaves without a question. It does not say which payments below the limit the attacker
can cause, or which other actions. That is the subject of the rest of this section.

\subsection{The Opaque Semantics}

An \emph{opaque value} is a value or the token $\op$. In an \emph{opaque expression}
variables have been replaced by opaque values, so $\op$ may occur where a value may. The
opaque value $\sem{t}$ of a term is $\op$ if $\op$ occurs in $t$, and its value otherwise.

The opaque semantics has a parameter $\kappa \in \{\E, \U\}$, the level from which data is
opaque. Figure~\ref{fig:opaque} defines the judgment $\oev{e}{\kappa}{o}{\tau}$ for an opaque
expression $e$, an opaque value $o$ and a trace $\tau$. An action whose result is untrusted
returns $\op$. The model returns $\op$ when it is asked about $\op$ (\textsc{O-Model}). A
request with $\op$ in its payload, a conditional with $\op$ in its test, and a check with
$\op$ among its parameters are not evaluated at all: they return $\op$ and leave no events
(\textsc{O-Skip}). The interesting rules are those for a check of a candidate that is $\op$.
At the level $\U$ the check makes a choice: \textsc{O-Choose} continues the first branch with
any value that satisfies the predicate, and \textsc{O-Refuse} continues with the second
branch, provided that some value fails the predicate. At the level $\E$ endorsed data is
opaque too, and the check is skipped.

\begin{figure}[t]
\begin{mathpar}
\inferrule*[right=O-Term]{ }{\oev{t}{\kappa}{\sem{t}}{\eps}}
\and
\inferrule*[right=O-Model]{o = \op \text{ if } \op \in \vec t,\ \text{else any value}}
  {\oev{\model{\vec t}}{\kappa}{o}{\eps}}
\and
\inferrule*[right=O-Act]{\op \notin \vec t \\ o = \op \text{ if } \ell_a = \U,\ \text{else any value}}
  {\oev{\act{a}{\vec t}}{\kappa}{o}{\kw{act}(a, \sem{\vec t})}}
\and
\inferrule*[right=O-Guard]{\op \notin \vec t \\ o = \op \text{ if } \ell_a = \U,\ \text{else any value}}
  {\oev{\grd{a}{\vec t}}{\kappa}{o}{\kw{grd}(a, \sem{\vec t})}}
\and
\inferrule*[right=O-Guard-Skip]{\op \in \vec t}{\oev{\grd{a}{\vec t}}{\E}{\op}{\eps}}
\and
\inferrule*[right=O-Yes]{\op \notin \vec t \\ o' = \op \text{ if } \ell_a = \U,\ \text{else any value} \\
  \oev{e_1\{o'/x\}}{\kappa}{o}{\tau}}
  {\oev{\ask{a}{\vec t}{x}{e_1}{e_2}}{\kappa}{o}{\kw{ok}(a, \sem{\vec t}) \cdot \tau}}
\and
\inferrule*[right=O-Skip]{e \text{ is } \kw{ask}\ldots \text{ with } \op \text{ in its payload, or }
  \kw{if}\ldots \text{ with } \op \text{ in its test, or} \\
  \kw{check}\ldots \text{ with } \op \text{ among its parameters, or with } \op
  \text{ in its candidate and } \kappa = \E}
  {\oev{e}{\kappa}{\op}{\eps}}
\and
\inferrule*[right=O-Choose]{\op \in t \\ \op \notin \varphi \\ \varphi\{c/x\} \text{ holds} \\
  \oev{e_1\{c/x\}}{\U}{o}{\tau}}
  {\oev{\chk{t}{x}{\varphi}{e_1}{e_2}}{\U}{o}{\tau}}
\and
\inferrule*[right=O-Refuse]{\op \in t \\ \op \notin \varphi \\ \varphi\{c/x\} \text{ fails for some } c \\
  \oev{e_2}{\U}{o}{\tau}}
  {\oev{\chk{t}{x}{\varphi}{e_1}{e_2}}{\U}{o}{\tau}}
\end{mathpar}
\Description{Nine inference rules of the opaque semantics.}
\caption{The opaque semantics at the level $\kappa$. The rules \textsc{O-No}, \textsc{O-Let},
\textsc{O-If}, \textsc{O-Pass}, \textsc{O-Fail}, \textsc{O-Iter0} and \textsc{O-Iter} are the
rules of Figure~\ref{fig:real} for expressions whose payload, test, candidate and parameters
contain no $\op$, with opaque values as results.}
\label{fig:opaque}
\end{figure}

The \emph{visible} part $\vis{\tau}{\kappa}$ of a trace consists of its unapproved actions
below the level $\kappa$: $\vis{\tau}{\U}$ keeps the events $\kw{act}$ and $\kw{grd}$, and
$\vis{\tau}{\E}$ keeps the events $\kw{act}$. We call $\vis{\tau}{\U}$ the \emph{clean
trace}. Approved actions are not visible at either level. They are vouched for by the person
who saw them, and the type system promises nothing more about them.

The opaque semantics puts no labels on data. What it treats as opaque is decided by the
values: the results of actions with $\ell_a = \U$, and everything computed from them.

\subsection{Every Attack Is a Sequence of Choices}

An opaque value $o$ \emph{matches} a value $c$ at the label $\ell$ if $o = c$, or $o = \op$
and $\kappa \sqsubseteq \ell$. An opaque substitution $\gamma^\circ$ matches $\gamma$ at
$\Gamma$ if $\gamma^\circ(x)$ matches $\gamma(x)$ at $\Gamma(x)$ for every $x$. So data below
the level $\kappa$ is the same in both, and data at or above $\kappa$ may be hidden.

\begin{lemma}[Quiet]
\label{lem:quiet}
If $\ty{\Phi}{pc}{e}{\ell\,\post{\psi}}$ and $\kappa \sqsubseteq pc$, then every evaluation
$\ev{e\gamma}{c}{\tau}$ has $\vis{\tau}{\kappa} = \eps$.
\end{lemma}

\begin{proof}
By induction on the derivation. \textsc{Act} requires $pc = \T$, which is below $\kappa$, so
it does not occur. \textsc{Act-$\pi$} requires $pc \sqsubseteq \E$, so $\kappa = \E$, and the
event $\kw{grd}$ is not visible at $\E$. All other rules type their premises at $pc$ or
above, and the events $\kw{ok}$ and $\kw{no}$ are never visible.
\end{proof}

\begin{lemma}[Totality]
\label{lem:total}
Every closed expression has an evaluation in the real semantics.
\end{lemma}

\begin{proof}
By induction on the expression, and on $n$ for a loop: terms and predicates are total.
\end{proof}

\begin{theorem}[Every attack is a sequence of choices]
\label{thm:abstraction}
Let $\ty{\Phi}{pc}{e}{\ell\,\post{\psi}}$, let $\kappa \in \{\E, \U\}$, and let
$\gamma^\circ$ match $\gamma$ at $\Gamma$. If $\ev{e\gamma}{c}{\tau}$, then there are $o$ and
$\tau^\circ$ such that
\[
\oev{e\gamma^\circ}{\kappa}{o}{\tau^\circ}, \qquad
\vis{\tau^\circ}{\kappa} = \vis{\tau}{\kappa}, \qquad
o \text{ matches } c \text{ at } \ell .
\]
\end{theorem}

\begin{proof}
By induction on the typing derivation. We use two facts about terms and predicates, which
follow from the definition of matching. For a term $t$, the opaque value of $t\gamma^\circ$
matches $\sem{t\gamma}$ at $\lab(t)$: either no variable of $t$ is hidden and the values are
equal, or some variable is hidden, its label is at least $\kappa$, and so is $\lab(t)$. For a
predicate $\varphi$, either $\varphi\gamma^\circ$ contains no $\op$ and holds exactly when
$\varphi\gamma$ holds, or it contains $\op$ and $\kappa \sqsubseteq \lab(\varphi)$.

\textsc{Term}: the first fact. \textsc{Model}: if $\vec t\gamma^\circ$ contains no $\op$,
\textsc{O-Model} returns the same value $c$; otherwise it returns $\op$, and
$\kappa \sqsubseteq \lab(\vec t)$.

\textsc{Act}: since $\lab(\vec t) = \T$ is below $\kappa$, the payloads are equal and contain
no $\op$. \textsc{O-Act} emits the same event, and it returns the same result $r$ if
$\ell_a = \T$ and $\op$ if $\ell_a = \U$; both match $r$ at $\ell_a$.

\textsc{Act-$\pi$}: if $\vec t\gamma^\circ$ contains no $\op$, \textsc{O-Guard} emits the
same event, and its result matches $r$ at $\ell_a \sqcup \lab(\vec t)$ as before. If it
contains $\op$, then $\kappa \sqsubseteq \lab(\vec t) \sqsubseteq \E$, so $\kappa = \E$.
\textsc{O-Guard-Skip} applies, the real event $\kw{grd}$ is not visible at $\E$, and $\op$
matches $r$ at a label that is at least $\E$.

\textsc{Ask}: if the payload $\vec t\gamma^\circ$ contains $\op$, then
$\kappa \sqsubseteq \ell$. \textsc{O-Skip} applies. The real evaluation continues with a
branch that is typed at $pc \sqcup \ell$, so its visible trace is empty by
Lemma~\ref{lem:quiet}, and the result label is at least $\ell$. Otherwise the payloads are
equal. If the real evaluation used \textsc{R-Yes} with the result $r$, we use \textsc{O-Yes}
with $o' = r$ if $\ell_a = \T$ and $o' = \op$ otherwise; $o'$ matches $r$ at
$\ell_a \sqcup \ell$, and the induction hypothesis applies to $e_1$ with the substitutions
extended by $x$. The case of \textsc{R-No} is similar.

\textsc{Let}: the induction hypothesis for $e_1$ gives $o_1$ that matches $c_1$ at $\ell_1$,
so $\gamma^\circ[x \mapsto o_1]$ matches $\gamma[x \mapsto c_1]$ at $\Gamma, x : \ell_1$, and
the hypothesis applies to $e_2$. The visible traces concatenate.

\textsc{If}: if $\varphi\gamma^\circ$ contains $\op$, then $\kappa \sqsubseteq \ell$.
\textsc{O-Skip} applies, the branch that the real evaluation runs is quiet by
Lemma~\ref{lem:quiet}, and the result label is at least $\ell$. Otherwise both evaluations
take the same branch, and the induction hypothesis applies.

\textsc{Check}: let $c_0 = \sem{t\gamma}$. There are four cases. (a) A parameter of the check
is hidden. Then $\kappa \sqsubseteq \lab(\varphi \setminus x) \sqsubseteq \ell_d$, and
\textsc{O-Skip} applies as in the case of \textsc{If}. (b) No parameter is hidden, the
candidate $t\gamma^\circ$ contains $\op$, and $\kappa = \E$. Then $\E \sqsubseteq \lab(t)$,
so $\ell_x = \E$ and $\kappa \sqsubseteq \ell_d$; again \textsc{O-Skip} applies. (c) No
parameter is hidden, the candidate contains $\op$, and $\kappa = \U$. If the real check
passed, then $\varphi\{c_0/x\}$ holds under $\gamma$, and it holds under $\gamma^\circ$ as
well because the parameters are equal. \textsc{O-Choose} with $c = c_0$ continues with
$e_1\{c_0/x\}$. The substitutions extended by $x \mapsto c_0$ match at
$\Gamma, x : \ell_x$, and the induction hypothesis applies to $e_1$. If the real check
failed, $c_0$ is a value that fails the predicate, \textsc{O-Refuse} applies, and the
hypothesis applies to $e_2$. (d) Nothing is hidden. Then the candidate and the outcome are
the same, \textsc{O-Pass} or \textsc{O-Fail} applies, and the hypothesis applies to the
branch.

\textsc{Iter}: by induction on $n$, with the induction hypothesis for the body at each turn.
\textsc{Sub}: a value that matches at $\ell$ matches at every label above $\ell$.
\end{proof}

For a closed program the theorem reads as follows. Let $e$ be closed with
$\cdot; \cdot; \T \vdash e : \ell\,\post{\psi}$. At the level $\U$: whatever the outside
world returns and whatever the model answers, the clean trace of the execution is the clean
trace of an opaque evaluation. That evaluation performs the same trusted and guarded
actions, although every value from outside is $\op$ to it. The only trace of the attacker in
it is the sequence of choices at \textsc{O-Choose} and \textsc{O-Refuse}.

\begin{corollary}[Trusted actions]
\label{cor:trusted}
For a closed well-typed program, the trusted actions of every real evaluation, in their
order, are the trusted actions of an opaque evaluation at the level $\E$. Such an evaluation
uses neither \textsc{O-Choose} nor \textsc{O-Refuse}.
\end{corollary}

This is the theorem at $\kappa = \E$. A trusted action is beyond the reach of the attacker
even through the choices that checks leave: an endorsed value cannot flow into it and cannot
decide whether it happens. It is the noninterference result of~\cite{tpp5}, extended to
programs with checks.

\subsection{Every Sequence of Choices Is an Attack}

The opaque semantics may offer more than an attacker can do. A check of a term such as
$f(u)$ for an untrusted $u$ can only see the values in the image of $f$, and a value that is
checked twice is the same value both times, while the opaque semantics chooses twice. These
differences disappear when checks are applied where guardrails usually sit, on what the
model has just produced. An expression is in \emph{generate-and-check form} if every check in
it occurs as $\letin{y}{\model{\vec t}}\; \chk{y}{x}{\varphi}{e_1}{e_2}$.

\begin{theorem}[Every sequence of choices is an attack]
\label{thm:exact}
Let $\ty{\Phi}{pc}{e}{\ell\,\post{\psi}}$ with $e$ in generate-and-check form, and let
$\gamma^\circ$ match $\gamma$ at $\Gamma$ for $\kappa = \U$. If
$\oev{e\gamma^\circ}{\U}{o}{\tau^\circ}$, then there are $c$ and $\tau$ such that
$\ev{e\gamma}{c}{\tau}$, $\vis{\tau}{\U} = \vis{\tau^\circ}{\U}$, and $o$ matches $c$ at
$\ell$.
\end{theorem}

\begin{proof}
By induction on the typing derivation, with the cases of the proof of
Theorem~\ref{thm:abstraction} read in the other direction. Where the opaque evaluation
returns a value from the model, the world or the person, the real evaluation returns the
same one; where it returns $\op$, the label of the result is $\U$, and the real evaluation
returns any value. Where the opaque evaluation skips an expression by \textsc{O-Skip}, the
real semantics has an evaluation of that expression by Lemma~\ref{lem:total}. The skipped
expression has $\op$ in its payload, its test or its parameters, so its branches are typed
at the program counter label $\U$ and are quiet by Lemma~\ref{lem:quiet}, and its own events
are not visible. \textsc{O-Guard-Skip} does not occur at the level $\U$.

The remaining case is a check of a candidate that is $\op$ with parameters that are not. By
the form of $e$, the check is the body of a $\kw{let}$ that binds its candidate $y$ to
$\model{\vec t}$, and $y$ is $\op$ because $\vec t\gamma^\circ$ contains $\op$. If the opaque
evaluation uses \textsc{O-Choose} with the value $c$, the real evaluation lets the model
answer $c$. The parameters are equal in both evaluations, so the real check passes, and both
continue with $e_1\{c/x\}$; the substitutions still match, with $y$ hidden. If the opaque
evaluation uses \textsc{O-Refuse}, some value $c'$ fails the predicate, the real evaluation
lets the model answer $c'$, and both continue with $e_2$.
\end{proof}

The proof lets the model answer what the attack needs. This is the assumption of the worst
case made concrete: a model that reads the attacker's text can be made to say anything. For
programs in generate-and-check form the two theorems give an equality: the set of clean
traces under all attacks is the set of clean traces of the opaque semantics.

\subsection{Counting}

The opaque semantics at the level $\U$ has two kinds of nondeterminism. The \emph{answers}
are the values that the model returns to questions without $\op$, the values that the world
returns to actions with trusted results, and the decisions of the person. The \emph{choices}
are the uses of \textsc{O-Choose}, with one option for each value that satisfies the
predicate, and of \textsc{O-Refuse}, which is one more option when some value fails it.
Suppose that the answers are given by functions of the request, so that only the choices
remain.

\begin{corollary}[Capacity]
\label{cor:capacity}
Let $e$ be closed and well typed, and let the answers be given by functions. Then the clean
trace of an opaque evaluation of $e$ at the level $\U$ is determined by its sequence of
choices. If every such evaluation makes at most $k$ choices, each among at most $s$ accepted
values, then $e$ has at most $(s+1)^k$ clean traces. More generally, the number of clean
traces is at most the largest product of the numbers of options along an evaluation.
\end{corollary}

\begin{proof}
With the answers fixed, every rule other than \textsc{O-Choose} and \textsc{O-Refuse}
determines its result and its premises, so the evaluations form a tree that branches only at
the choices. The number of leaves of a tree is at most the largest product of the branching
degrees along a path, by induction on the tree.
\end{proof}

By Theorem~\ref{thm:abstraction}, these are all the clean traces that an attacker can cause,
against a model, a world and a person that answer by the given functions wherever the
attacker's data is not involved. In bits, a check with $s$ accepted values gives the
attacker at most $\log_2(s+1)$ bits of control over the unapproved actions each time it
runs.

\begin{corollary}[No check, no influence]
\label{cor:nocheck}
Let $e$ be closed, well typed, and without checks, and let the answers be given by
functions. Then $e$ has exactly one clean trace.
\end{corollary}

\begin{proof}
Without checks there are no choices, so Corollary~\ref{cor:capacity} gives at most one
clean trace. There is one: by Lemma~\ref{lem:total} the program has a real evaluation whose
answers are given by the functions, and the opaque evaluation that
Theorem~\ref{thm:abstraction} constructs from it uses the same answers.
\end{proof}

\section{Six Laws for Guardrails}
\label{sec:laws}

This section reads the theorems as advice. Each law comes with a program, and the script of
Section~\ref{sec:tests} confirms the verdict of the type checker on each. In the script the
value set is $\{0, \ldots, 4\}$ and the limit $L$ of payments is $3$, so its numbers are
small; the programs are those shown here.

\subsection{A Model Cannot Endorse}
\label{sec:judge}

The program of Section~\ref{sec:overview} that asks a model to judge a payment is rejected,
and the reason is not particular to that program. By Corollary~\ref{cor:nocheck}, a program
without checks has a single clean trace: for given answers to its trusted questions, its
unapproved actions are the same under every attack. A program may call the model on
untrusted data as often as it likes, to classify it, to summarize it, or to look for
attacks in it. Whatever comes back is untrusted, and the type system keeps it away from
every action that is not approved. Endorsement is the privilege of code and of people.

This is a statement about guarantees and says nothing against the practical value of such
classifiers. A classifier that catches most injections lowers the rate of attacks that
reach a check or a person. It does not shrink the set of things that can happen, which is
what this paper measures.

\subsection{Reference Data Must Be Trusted}
\label{sec:reference}

A check compares a candidate with something: a limit, a list of contacts, a pattern. Where
that something comes from matters as much as the comparison.
\begin{prog}
\letin{l}{\model{m}} & l : \U \text{, a limit taken from the message} \\
\chk{p}{x}{x \le l}{\grd{\nm{pay}}{x}}{0} & \ell_d = \U \text{: rejected}
\end{prog}
In \textsc{Check} the label of the decision includes the labels of the parameters. With an
untrusted limit, $\ell_d = \U$, the first branch runs at an untrusted program counter, and
\textsc{Act-$\pi$} rejects the payment. The opaque semantics shows why this is right: with
$\op$ among its parameters the check is skipped, because an attacker who controls the limit
controls the outcome for every candidate. The same program with the limit read from the
user is accepted, provided the facts entail the policy:
\begin{prog}
\letin{l}{\act{\nm{read}}{}} & l : \T \\
\kw{check}\;p\;\kw{as}\;x.\,x \le l\;\kw{then} & x : \E \text{, fact } x \le l \\
\quad \ite{l \le L}{\grd{\nm{pay}}{x}}{\ldots} & x \le l,\; l \le L \models x \le L
\end{prog}
A common instance is an allowlist of recipients that the agent has itself fetched from a
place the attacker can write to.

\subsection{A Check Is Worth Its Acceptance Set}
\label{sec:bits}

How much does a check leave to the attacker? Corollary~\ref{cor:capacity} answers with a
number. An agent that lets the model classify a message into one of three queues and files
it accordingly,
\begin{prog}
\letin{k}{\model{m}} \\
\kw{check}\;k\;\kw{as}\;x.\,1 \le x \wedge x \le 3 \\
\quad \kw{then}\;\ite{x = 1}{\grd{\nm{file}}{1}}{\ite{x = 2}{\grd{\nm{file}}{2}}{\grd{\nm{file}}{3}}} \\
\quad \kw{else}\;0
\end{prog}
has exactly \exRoutes{} clean traces: the three queues and no action. The attacker controls
two bits. The payment agent of Section~\ref{sec:overview} has $L + 2$ clean traces, which
the script confirms as \exInvoiceTraces{} for $L = 3$. A check that validates free text
against a format has an acceptance set that is astronomically large, and the corollary then
says that the check is worth little: the attacker chooses among all well-formed texts. A
guardrail is strong in proportion to how few values it lets through, and the theorem makes
the count a quantity that can be compared across designs. By Theorem~\ref{thm:exact} the
count is exact for programs in generate-and-check form.

\subsection{Limits Do Not Add Up, Budgets Do}
\label{sec:budget}

A policy constrains one action. An agent that handles $n$ invoices in a loop, each with the
check $x \le L$, satisfies the policy at every payment and can still be made to pay $nL$ in
total: the opaque semantics chooses $L$ at every turn, and by Theorem~\ref{thm:exact} an
attack realizes the choice. The remedy is to make the total a part of the state and of the
check:
\begin{prog}
\kw{iter}\;n\;\kw{from}\;0\;\kw{as}\;s.\ & \text{invariant } \nu \le B \\
\quad \letin{m}{\act{\nm{fetch}}{\nm{inbox}}}\; \letin{p}{\model{m}} & m, p : \U \\
\quad \kw{check}\;p\;\kw{as}\;x.\,x \le L \wedge s + x \le B & x : \E \\
\quad\quad \kw{then}\;\letin{r}{\grd{\nm{pay}}{x}}\; s + x & \nu = s + x \text{ and } s + x \le B \\
\quad\quad \kw{else}\;s & \nu = s \text{ and } s \le B
\end{prog}
The state $s$ is the sum of the payments so far. It has the label $\E$, since endorsed
amounts flow into it, and that is allowed for a parameter of a check: the decision of the
check then has the label $\E$. The invariant $\nu \le B$ holds at the start, and each branch
of the check preserves it, by the fact that the check records and by the invariant on $s$.
By Theorem~\ref{thm:policy} the result of the loop is at most $B$, and the result is the
total paid. The attacker now chooses a sequence of payments with a sum of at most $B$.

In the script, with $n = 4$, $L = 3$ and $B = 3$, the largest total over all clean traces is
\exTotalWithout{} without the budget and \exTotalWith{} with it, and the number of clean
traces falls from \exTracesWithout{} to \exTracesWith. A budget is a check whose parameter is
endorsed state. The type system admits such parameters and keeps the bookkeeping honest: if
the program forgot to add a payment to $s$, the invariant would still hold and the theorem
would still bound $s$, but $s$ would no longer be the total, and that is a property of the
program for which a stronger postcondition is needed.

\subsection{The Result of a Guarded Action Is Endorsed}
\label{sec:result}

The world answers to what it was sent. If an attacker chose the payload of an action within
its policy, the result of the action is under the attacker's influence even when the action
has a trusted result label. \textsc{Act-$\pi$} therefore gives the result the label
$\ell_a \sqcup \lab(\vec t)$. With $k$ an untrusted answer of the model as before, the
following program uses the result of a guarded action in a trusted action and is rejected:
\begin{prog}
\letin{y}{\chk{k}{x}{x \le 4}{x}{0}} & y : \E \\
\letin{r}{\grd{\nm{file}}{y}} & r : \E \text{, not } \T \\
\act{\nm{log}}{r} & \text{rejected}
\end{prog}
We did not see the need for this join when we first wrote the rule. The mutation test of
Section~\ref{sec:tests} breaks each rule in turn, and a checker whose guarded actions
return results of label $\ell_a$ alone accepts this program, on which the opaque semantics
at the level $\E$ is stuck: it would have to perform a trusted action on $\op$. A lookup in
a trusted database under a key that the attacker selected is the everyday form of this
flow.

\subsection{The Attacker Can Always Refuse}
\label{sec:refuse}

The rule \textsc{O-Refuse} is as much a part of the attacker's power as \textsc{O-Choose}:
an attacker can always make the model produce a value that fails the check. Whatever a
program does in the second branch of a check is therefore the attacker's to trigger. The
payment agent of Section~\ref{sec:overview} falls back to asking the person. In a loop over
$n$ messages the attacker can cause $n$ questions, each about a payload of the attacker's
choosing. This turns a guardrail into a way of spending the attention of the person, which
is the resource that selective approval was meant to protect~\cite{tpp5}. The opaque
semantics does not list these questions among the clean traces, since approved actions are
outside its guarantee, but it shows the choice that leads to them. A second branch that
abstains, or a counter in the state of the loop that bounds the number of questions as the
budget bounds the payments, removes the lever.

The same reading applies to availability. Refusal is always possible, so no check-guarded
action can be guaranteed to happen under attack. Checks buy integrity and sell liveness,
and an earlier paper in this series prices that trade for an honest model that merely
errs~\cite{tpp3}.

\section{Testing the Claims}
\label{sec:tests}

The artifact is one Python script, \texttt{flow.py}, of about 840 lines without
dependencies. It instantiates the set of values with $\{0, \ldots, 4\}$, so that entailment
between predicates is decided by enumeration, and it contains the following parts.

\paragraph{A type checker.} It implements Figure~\ref{fig:typing} in an algorithmic form
that computes the least label of each expression and checks postconditions against a given
predicate. The invariants of loops are written in the programs.

\paragraph{Two interpreters.} One implements Figure~\ref{fig:real} and one
Figure~\ref{fig:opaque}. Each returns the value, the events and a log of its decisions, and
each can be \emph{guided} by a log of the other: the opaque interpreter then takes the
answers and the outcomes of the checks from a real evaluation and fails if a rule of the
opaque semantics does not permit the step, and the real interpreter takes the answers and
the choices from an opaque evaluation and lets the model produce the chosen values. Guided
evaluation is the construction in the proofs of Theorems~\ref{thm:abstraction}
and~\ref{thm:exact}, as a program.

\paragraph{Random programs.} One generator builds programs that are well typed by
construction. A second one builds programs without regard to the rules, and the tests keep
those that the checker accepts. For \exPrograms{} programs and four real evaluations of
each, the script checks that every guarded event satisfies its policy
(Theorem~\ref{thm:policy}) and that the guided opaque evaluation exists at both levels, has
the same visible trace, and returns a matching value (Theorem~\ref{thm:abstraction}). For
\exPaired{} programs in generate-and-check form and four opaque evaluations of each with
random choices, it checks that the guided real evaluation exists and has the same clean
trace (Theorem~\ref{thm:exact}). All checks pass.

\paragraph{Broken rules.} Passing tests show little unless they can fail. The checker has a
switch that breaks one rule at a time: the answer of the model is trusted, a trusted action
accepts any payload, a guarded action accepts untrusted influence, the policy is not
checked, the result of a guarded action has the label $\ell_a$ only, the branches of
\textsc{Ask}, \textsc{If} or \textsc{Check} do not raise the program counter, the results of
these three do not include the label of the decision, a $\kw{let}$ forgets the label of the
bound expression, a check endorses to $\T$, a check ignores the labels of its parameters,
and a loop does not iterate its label to a fixed point. For each of these \exRules{} broken
checkers the script searches for a program that the broken checker accepts and on which a
test fails, and it finds one in every case. The search tries the program of
Section~\ref{sec:result} first, which exposes five of the broken rules, and then random
programs. The slowest case, a conditional whose result forgets the label of its test, needs
\exSlowest{} candidates. We added the program of Section~\ref{sec:result} to the search
after random programs alone had failed to expose the missing join in the result of a guarded
action; they rarely contain that flow.

The mutation test also found a label that carries no weight. In \textsc{Ask} the variable
$x$ could have the label $\ell_a$ alone: the program counter of its branch already includes
the label of the payload, and the proof of Theorem~\ref{thm:abstraction} goes through
unchanged. No test rejects that variant. We keep the label $\ell_a \sqcup \ell$, which says
what the result depends on.

\paragraph{The examples.} The script runs the checker on the programs of
Sections~\ref{sec:overview} and~\ref{sec:laws} and confirms each verdict. It enumerates all
sequences of choices of the classification agent, the payment agent and the two loops of
Section~\ref{sec:budget}, with answers given by fixed functions, and reports the numbers of
clean traces quoted above together with the bound of Corollary~\ref{cor:capacity}; for the
loop with a budget there are \exTracesWith{} clean traces against a bound of \exBoundWith.
The command \texttt{make -C artifact check} runs everything in about ten seconds and
compares the output with the committed copy.

\section{Related Work}
\label{sec:related}

\paragraph{Integrity and information flow.} Labels on data and a program counter label go
back to the lattice model of Denning~\cite{denning1976} and to type systems for secure
information flow~\cite{volpano1996,sabelfeld2003}. {\O}rb{\ae}k and
Palsberg~\cite{orbaek1997} analyze trust in a $\lambda$-calculus with constructs that mark
data as trusted or distrusted and a construct that requires trust. Our labels and our rules
for conditionals are in this tradition.

\paragraph{Endorsement and its semantics.} Noninterference forbids all influence, and
practical systems need controlled exceptions. For confidentiality the exception is
declassification, whose dimensions Sabelfeld and Sands~\cite{sabelfeld2009} classify, and
for integrity it is endorsement. Robust declassification~\cite{zdancewic2001} asks that an
attacker cannot influence what is released, and qualified
robustness~\cite{myers2006} extends it to programs that endorse. Askarov and
Myers~\cite{askarov2011} characterize security by the attacker's influence, define the
attacker's impact on trusted locations, and give conditions and a type system for
\emph{checked endorsement}, in which a value is endorsed only if a test on it succeeds.
The check of our calculus is a checked endorsement in their sense, and the idea of stating
the guarantee as a bound on the attacker's influence is theirs. Our formulation differs in
three respects. The guarantee is a refinement between two semantics of the same program,
one of which never inspects untrusted data, so that the attacked program can be verified as
an ordinary nondeterministic program. The refinement is exact for a natural class of
programs. And the endorsed value carries the predicate of its check as a refinement, which
connects endorsement to the policies of actions. Relaxed
noninterference~\cite{li2005} specifies downgrading by the programs through which data may
pass; a check with its predicate is such a specification for integrity. Clarkson and
Schneider~\cite{clarkson2010} quantify integrity with information theory, as the
contamination of trusted data by untrusted inputs. Our capacity bound is a possibilistic
count in the same spirit and uses no probabilities.

\paragraph{Refinements.} Recording the predicate of a test as a fact about a variable is the
basic move of refinement types~\cite{freeman1991,rondon2008}. We use a deliberately small
fragment, with entailment as an oracle.

\paragraph{Prompt injection.} Greshake et al.~\cite{greshake2023} showed that text retrieved
by an application can take control of the model in it. CaMeL~\cite{debenedetti2025} extracts
the control flow and the data flow from the trusted request of the user, so that untrusted
data cannot affect the flow of the program, and enforces policies with capabilities when
tools are called. Costa et al.~\cite{costa2025} give a formal model of agent planners,
characterize what dynamic taint tracking can enforce, and build a planner that tracks
integrity and confidentiality labels. Beurer-Kellner et al.~\cite{beurerkellner2025} propose
design patterns for agents with provable resistance to injection and discuss their cost in
utility, and AgentDojo~\cite{debenedetti2024} measures attacks and defenses on agents with
tools. Safeguard models such as Llama Guard~\cite{inan2023} classify inputs and outputs with
a model; Section~\ref{sec:judge} says what such a classifier can and cannot contribute to a
guarantee. These systems track flows at run time or fix an architecture. We give a static
discipline, and a theorem that says what remains possible once checks have been placed:
which is what a designer must know to decide whether a pattern is strong enough.

\paragraph{This series.} The two-level system without endorsement is~\cite{tpp5}, and its
approvals are the linear tokens of~\cite{tpp4}. The same construct, a check of what a model
generated, is analyzed for honest models in~\cite{tpp3}: there a check turns errors into
abstentions with a probability, and here it turns an attack into a choice within a set.

\section{Limits and Open Problems}
\label{sec:limits}

\paragraph{Limits.} The attacker of this paper is as strong as possible, and the results
are possibilistic: they say which traces can occur and do not say how likely an attack is to
steer a real model to a given choice. The labels $\ell_a$ are declarations. If an approved
action with an untrusted payload changes the state of the world, and a later action with a
trusted result reads that state, the declaration is wrong, and no theorem here notices. The
person is modelled as nondeterminism; a person who approves without reading provides no
protection, and the calculus cannot tell. Predicates are total functions on abstract values.
For real payloads, which are strings and structured documents, the difficulty is to write
predicates whose acceptance sets are small and still useful. The calculus has bounded loops
and no functions. The proofs are not mechanized.

\paragraph{Open problems.}
\begin{enumerate}
\item \emph{Probabilities.} With a distribution on the answers of an honest model and an
attacker who controls only a part of the input, the capacity of a check becomes a quantity
of information, and it should compose with the reliability grades of~\cite{tpp3}. A single
judgment could then bound both the error of an agent that is not attacked and the influence
on one that is.
\item \emph{Confidentiality.} The dual system has secret data, public sinks, and checks that
declassify. An agent that reads a secret and writes to an address needs both systems at
once, and the known interaction between them is robustness: the attacker must not choose
what is declassified.
\item \emph{Choosing where to check.} Given a program without checks and policies for its
actions, synthesize the checks with the smallest total capacity that make it well typed, and
ask a person for the rest.
\item \emph{Structured values.} For strings, an acceptance set is a language. Its size at a
given length, or the entropy of the checker's grammar, would be the measure that replaces
$\log_2(s+1)$.
\item \emph{Monitors.} The static discipline has a dynamic counterpart with labels at run
time, which would let a harness enforce it on programs that a model writes. The opaque
semantics suggests a design: run the planner on $\op$.
\item \emph{Mechanization.} Theorems~\ref{thm:abstraction} and~\ref{thm:exact} are
inductions over a first-order big-step semantics, and the guided interpreters of the
artifact are their computational content.
\end{enumerate}

\section{Conclusion}

A guardrail that is ordinary code turns an injection into a choice. The attacker picks one
of the values that the check accepts, or makes the check fail, and has no other influence
on what the agent does without approval. The type system finds the places where a program
relies on more than that, and the opaque semantics is the program that remains: no model, no
untrusted text, and a choice wherever a check stands. Guardrails can then be compared by
counting, and agents can be verified against their opaque semantics by the means that
already exist for nondeterministic programs.

\bibliographystyle{ACM-Reference-Format}
\bibliography{refs}

\end{document}
